% GENERATED FILE. Edit canonical lessons, then run npm run generate:books.
% canonical-source-hash: fnv1a64:6032b72ff653c6a1
% canonical-lessons: SA-C546-santustah, SA-C546-nispannah, SA-C546-udbhavati, SA-C546-praptih, SA-C546-vipakah

\chapter{How It Turned Out}
\label{ch:sa-how-it-turned-out}
\hlchaptermodality{546}

\begin{chapteropening}
\textbf{By the end of this chapter:} \emph{I can say how something turned out.}

Complete the last lesson of chapter 546: I can say how something turned out.
\end{chapteropening}

\section[\texorpdfstring{\sk{सन्तुष्टः} (santuṣṭaḥ)}{सन्तुष्टः (santuṣṭaḥ)}]{\sk{सन्तुष्टः} (santuṣṭaḥ) — satisfied; contented}

\label{lesson:SA-C546-santustah}



\begin{quote}
Before the new one: say the Sanskrit for at first, then the Sanskrit for reviews.
\end{quote}

\subsection*{You'll want to know: \sk{सन्तुष्टः}}
\textbf{\sk{सन्तुष्टः}} — \emph{santuṣṭaḥ} — \textquotedblleft{}satisfied; contented\textquotedblright{}.

\begin{quote}
\sk{यदि} \sk{परिणामः} \sk{शुभः} \sk{अस्ति}, \sk{तर्हि} \sk{पिता} \sk{सन्तुष्टः} \sk{भवति।}
\end{quote}

\begin{quote}
\emph{yadi pariṇāmaḥ śubhaḥ asti, tarhi pitā santuṣṭaḥ bhavati.}
\end{quote}

\begin{quote}
If the result is good, father is satisfied.
\end{quote}

\begin{cousinweb}
\textbf{\sk{सन्तुष्टः}} (\emph{santuṣṭaḥ}) is built on \emph{sam}, \textquotedblleft{}fully\textquotedblright{}, and \emph{tuṣṭa}, \textquotedblleft{}pleased\textquotedblright{}, the past participle of the root √tuṣ, \textquotedblleft{}to be pleased\textquotedblright{}, which also gives \emph{tuṣyati}; MW gives \emph{saṃtuṣṭa} as quite satisfied, contented, well pleased (MW).
\end{cousinweb}

\begin{grammarlens}[title={What you've built}]
One more word for saying what happens if.
\end{grammarlens}

\subsection*{Your turn}
\begin{itemize}
\raggedright
  \item \emph{Say it:} \emph{santuṣṭaḥ}
  \item \emph{Say it:} \emph{santuṣṭaḥ}, once more
  \item \emph{Say it:} \emph{yadi pariṇāmaḥ śubhaḥ asti, tarhi pitā santuṣṭaḥ bhavati.} — If the result is good, father is satisfied.
  \item \emph{Recall:} say the Sanskrit for at first, then the Sanskrit for reviews, then say \emph{santuṣṭaḥ} again
\end{itemize}

\begin{tcolorbox}[breakable,colback=teal!4,colframe=teal!35!black,title={Before you move on}]
What does \sk{सन्तुष्टः} mean? (\textquotedblleft{}Satisfied; contented\textquotedblright{}.) Say it once more.
\end{tcolorbox}
\section[\texorpdfstring{\sk{निष्पन्नः} (niṣpannaḥ)}{निष्पन्नः (niṣpannaḥ)}]{\sk{निष्पन्नः} (niṣpannaḥ) — completed}

\label{lesson:SA-C546-nispannah}



\begin{quote}
Before the new one: say the Sanskrit for reviews, then the Sanskrit for satisfied.
\end{quote}

\subsection*{You'll want to know: \sk{निष्पन्नः}}
\textbf{\sk{निष्पन्नः}} — \emph{niṣpannaḥ} — \textquotedblleft{}completed\textquotedblright{}.

\begin{quote}
\sk{यदि} \sk{कार्यं} \sk{निष्पन्नम्} \sk{अस्ति}, \sk{तर्हि} \sk{वयं} \sk{गृहं} \sk{गच्छामः।}
\end{quote}

\begin{quote}
\emph{yadi kāryaṃ niṣpannam asti, tarhi vayaṃ gṛhaṃ gacchāmaḥ.}
\end{quote}

\begin{quote}
If the work is done, we go home.
\end{quote}

\begin{cousinweb}
\textbf{\sk{निष्पन्नः}} (\emph{niṣpannaḥ}) is built on \emph{nis}, \textquotedblleft{}out\textquotedblright{}, and the root √pad, \textquotedblleft{}to go\textquotedblright{}, as its past participle; MW gives \emph{niṣpanna} as brought about, completed, finished, ready (MW).
\end{cousinweb}

\begin{grammarlens}[title={What you've built}]
One more word for saying what happens if.
\end{grammarlens}

\subsection*{Your turn}
\begin{itemize}
\raggedright
  \item \emph{Say it:} \emph{niṣpannaḥ}
  \item \emph{Say it:} \emph{niṣpannaḥ}, once more
  \item \emph{Say it:} \emph{yadi kāryaṃ niṣpannam asti, tarhi vayaṃ gṛhaṃ gacchāmaḥ.} — If the work is done, we go home.
  \item \emph{Recall:} say the Sanskrit for reviews, then the Sanskrit for satisfied, then say \emph{niṣpannaḥ} again
\end{itemize}

\begin{tcolorbox}[breakable,colback=teal!4,colframe=teal!35!black,title={Before you move on}]
What does \sk{निष्पन्नः} mean? (\textquotedblleft{}Completed\textquotedblright{}.) Say it once more.
\end{tcolorbox}
\section[\texorpdfstring{\sk{उद्भवति} (udbhavati)}{उद्भवति (udbhavati)}]{\sk{उद्भवति} (udbhavati) — arises}

\label{lesson:SA-C546-udbhavati}



\begin{quote}
Before the new one: say the Sanskrit for satisfied, then the Sanskrit for completed.
\end{quote}

\subsection*{You'll want to know: \sk{उद्भवति}}
\textbf{\sk{उद्भवति}} — \emph{udbhavati} — \textquotedblleft{}arises\textquotedblright{}.

\begin{quote}
\sk{यदि} \sk{प्रश्नः} \sk{उद्भवति}, \sk{तर्हि} \sk{अध्यापकं} \sk{पृच्छतु।}
\end{quote}

\begin{quote}
\emph{yadi praśnaḥ udbhavati, tarhi adhyāpakaṃ pṛcchatu.}
\end{quote}

\begin{quote}
If a question comes up, ask the teacher.
\end{quote}

Past (it …): \sk{प्रश्नः} \sk{उद्भूतवान्}, \emph{praśnaḥ udbhūtavān}.

\begin{cousinweb}
\textbf{\sk{उद्भवति}} (\emph{udbhavati}) is built on \emph{ud}, \textquotedblleft{}up\textquotedblright{}, and the root √bhū, \textquotedblleft{}to become\textquotedblright{}, which also gives \emph{bhavati}; MW gives \emph{udbhū} as to come forth, to arise, to spring from; its past participle is \emph{udbhūta} (MW).
\end{cousinweb}

\begin{grammarlens}[title={What you've built}]
One more word for saying what happens if.
\end{grammarlens}

\subsection*{Your turn}
\begin{itemize}
\raggedright
  \item \emph{Say it:} \emph{udbhavati}
  \item \emph{Say it:} \emph{udbhavati}, once more
  \item \emph{Say it:} \emph{yadi praśnaḥ udbhavati, tarhi adhyāpakaṃ pṛcchatu.} — If a question comes up, ask the teacher.
  \item \emph{Recall:} say the Sanskrit for satisfied, then the Sanskrit for completed, then say \emph{udbhavati} again
\end{itemize}

\begin{tcolorbox}[breakable,colback=teal!4,colframe=teal!35!black,title={Before you move on}]
What does \sk{उद्भवति} mean? (\textquotedblleft{}Arises\textquotedblright{}.) Say it once more.
\end{tcolorbox}
\section[\texorpdfstring{\sk{प्राप्तिः} (prāptiḥ)}{प्राप्तिः (prāptiḥ)}]{\sk{प्राप्तिः} (prāptiḥ) — an attainment}

\label{lesson:SA-C546-praptih}



\begin{quote}
Before the new one: say the Sanskrit for completed, then the Sanskrit for arises.
\end{quote}

\subsection*{You'll want to know: \sk{प्राप्तिः}}
\textbf{\sk{प्राप्तिः}} — \emph{prāptiḥ} — \textquotedblleft{}an attainment\textquotedblright{}.

\begin{quote}
\sk{यदि} \sk{लक्ष्यस्य} \sk{प्राप्तिः} \sk{भवति}, \sk{तर्हि} \sk{वयम्} \sk{उत्सवं} \sk{करिष्यामः।}
\end{quote}

\begin{quote}
\emph{yadi lakṣyasya prāptiḥ bhavati, tarhi vayam utsavaṃ kariṣyāmaḥ.}
\end{quote}

\begin{quote}
If we reach our goal, we will celebrate.
\end{quote}

A feminine noun.

\begin{cousinweb}
\textbf{\sk{प्राप्तिः}} (\emph{prāptiḥ}) is built on \emph{pra}, \textquotedblleft{}forth\textquotedblright{}, and the root √āp, \textquotedblleft{}to reach\textquotedblright{}, which also gives \emph{prāpnoti}; MW gives \emph{prāpti} as reaching, obtaining, attainment (MW).
\end{cousinweb}

\begin{grammarlens}[title={What you've built}]
One more word for saying what happens if.
\end{grammarlens}

\subsection*{Your turn}
\begin{itemize}
\raggedright
  \item \emph{Say it:} \emph{prāptiḥ}
  \item \emph{Say it:} \emph{prāptiḥ}, once more
  \item \emph{Say it:} \emph{yadi lakṣyasya prāptiḥ bhavati, tarhi vayam utsavaṃ kariṣyāmaḥ.} — If we reach our goal, we will celebrate.
  \item \emph{Recall:} say the Sanskrit for completed, then the Sanskrit for arises, then say \emph{prāptiḥ} again
\end{itemize}

\begin{tcolorbox}[breakable,colback=teal!4,colframe=teal!35!black,title={Before you move on}]
What does \sk{प्राप्तिः} mean? (\textquotedblleft{}An attainment\textquotedblright{}.) Say it once more.
\end{tcolorbox}
\section[\texorpdfstring{\sk{विपाकः} (vipākaḥ)}{विपाकः (vipākaḥ)}]{\sk{विपाकः} (vipākaḥ) — a consequence}

\label{lesson:SA-C546-vipakah}



\begin{quote}
Before the new one: say the Sanskrit for arises, then the Sanskrit for an attainment.
\end{quote}

\subsection*{You'll want to know: \sk{विपाकः}}
\textbf{\sk{विपाकः}} — \emph{vipākaḥ} — \textquotedblleft{}a consequence\textquotedblright{}.

\begin{quote}
\sk{यदि} \sk{कर्म} \sk{उत्तमम्} \sk{अस्ति}, \sk{तर्हि} \sk{विपाकः} \sk{अपि} \sk{उत्तमः।}
\end{quote}

\begin{quote}
\emph{yadi karma uttamam asti, tarhi vipākaḥ api uttamaḥ.}
\end{quote}

\begin{quote}
If a deed is good, its outcome is good too.
\end{quote}

A masculine noun.

\begin{cousinweb}
\textbf{\sk{विपाकः}} (\emph{vipākaḥ}) is built on \emph{vi} and the root √pac, \textquotedblleft{}to cook, to ripen\textquotedblright{}, which also gives \emph{pacati}; MW gives \emph{vipāka} as ripening, maturing, the result or consequence of actions (MW).
\end{cousinweb}

\begin{grammarlens}[title={What you've built}]
One more word for saying what happens if.
\end{grammarlens}

\subsection*{Your turn}
\begin{itemize}
\raggedright
  \item \emph{Say it:} \emph{vipākaḥ}
  \item \emph{Say it:} \emph{vipākaḥ}, once more
  \item \emph{Say it:} \emph{yadi karma uttamam asti, tarhi vipākaḥ api uttamaḥ.} — If a deed is good, its outcome is good too.
  \item \emph{Recall:} say the Sanskrit for arises, then the Sanskrit for an attainment, then say \emph{vipākaḥ} again
\end{itemize}

\begin{tcolorbox}[breakable,colback=teal!4,colframe=teal!35!black,title={Before you move on}]
What does \sk{विपाकः} mean? (\textquotedblleft{}A consequence\textquotedblright{}.) Say it once more.
\end{tcolorbox}
